\documentclass[11pt,spanish]{article}
\renewcommand{\sfdefault}{lmss}
\renewcommand{\familydefault}{\rmdefault}
\usepackage[T1]{fontenc}
\usepackage[utf8]{inputenc}
\usepackage{babel}
\addto\shorthandsspanish{\spanishdeactivate{~<>.}}
\deactivatequoting
\usepackage{geometry}
\geometry{tmargin=2cm,bmargin=2cm,lmargin=2cm,rmargin=2cm}
\usepackage{fancyhdr}
\pagestyle{fancy}
\setlength{\headheight}{14pt}
\usepackage[skip=\bigskipamount]{parskip}
\usepackage{array}
\usepackage{longtable}
\usepackage{booktabs}
\usepackage{microtype}
\usepackage{hyperref}
\hypersetup{hidelinks}
\usepackage{pifont}
\usepackage{titling}
\setlength{\droptitle}{-2cm}
\usepackage{xcolor}
\usepackage{enumitem}
\usepackage[most]{tcolorbox}
\usepackage{amssymb}
\usepackage{amsmath}
\usepackage{tikz}
\usetikzlibrary{arrows.meta,automata,positioning}
\usepackage{listings}

\definecolor{darkpastelblue}{rgb}{0.47,0.62,0.8}
\definecolor{amber}{rgb}{1.0,0.75,0.0}
\definecolor{rojo7a3c3d}{rgb}{0.478,0.235,0.239}
\definecolor{verde587246}{rgb}{0.345,0.447,0.275}
\definecolor{antiquefuchsia}{rgb}{0.57,0.36,0.51}
\definecolor{pastelgreen}{rgb}{0.82,0.92,0.78}
\definecolor{pastelblue}{rgb}{0.78,0.87,0.96}
\definecolor{pastelpeach}{rgb}{0.98,0.84,0.72}
\definecolor{indiagreen}{rgb}{0.07,0.53,0.03}
\definecolor{codegray}{rgb}{0.96,0.96,0.96}

\newcommand{\blank}[1][4cm]{\rule{#1}{0.4pt}}
\newcommand{\cuadro}{\(\square\)}
\newcommand{\cf}[2]{\colorbox{#1}{\textcolor{white}{\strut #2}}}
\newcommand{\entrada}{\textbf{Entrada}}
\newcommand{\transformacion}{\textbf{Transformación}}
\newcommand{\salida}{\textbf{Salida}}

\newtcolorbox{estacion}[1]{
  colback=antiquefuchsia!5!white,
  colframe=antiquefuchsia!85!black,
  fonttitle=\bfseries,
  title={#1},
  breakable}

\newtcolorbox{modelo}[1]{
  colback=darkpastelblue!10!white,
  colframe=darkpastelblue!85!black,
  fonttitle=\bfseries,
  title={#1},
  breakable}

\newtcolorbox{checkpoint}[1]{
  colback=amber!7!white,
  colframe=amber!85!black,
  fonttitle=\bfseries,
  title={#1},
  breakable}

\newtcolorbox{evidencia}[1]{
  colback=pastelgreen!35!white,
  colframe=indiagreen!85!black,
  fonttitle=\bfseries,
  title={#1},
  breakable}

\newtcolorbox{decision}[1]{
  colback=pastelpeach!28!white,
  colframe=rojo7a3c3d!85!black,
  fonttitle=\bfseries,
  title={#1},
  breakable}

\lstdefinestyle{pds}{
  language=Python,
  basicstyle=\ttfamily\small,
  backgroundcolor=\color{codegray},
  frame=single,
  rulecolor=\color{black!30},
  breaklines=true,
  showstringspaces=false,
  columns=fullflexible,
  keepspaces=true,
  keywordstyle=\bfseries\color{indiagreen!70!black},
  commentstyle=\itshape\color{rojo7a3c3d},
  stringstyle=\color{antiquefuchsia!85!black}
}
\lstset{style=pds}

\setlength{\parindent}{0pt}
\setlist[itemize]{leftmargin=*,itemsep=0.35em}
\setlist[enumerate]{leftmargin=*,itemsep=0.45em}
\renewcommand{\arraystretch}{1.25}

\begin{document}

\lhead{PDS}
\rhead{Práctica 01}
\title{Programación de Sistemas\\
Temas 1.1 a 2.3\\ \bigskip
\cf{rojo7a3c3d}{\textbf{Práctica 1}}\\
\bigskip
\Large\textbf{Del programa fuente a los tokens:\\analizador léxico para \textsc{IMP}}}
\author{Manuel Soto Romero \and Araceli Liliana Reyes Cabello}
\date{Semestre 2026-2 \\ Colegio de Ciencia y Tecnología, UACM}

\maketitle

Esta práctica acompaña la transición desde el panorama de la programación de sistemas hasta la construcción del primer componente del compilador. Avanzaremos, de manera guiada, hacia las expresiones regulares, los autómatas finitos y la implementación de un analizador léxico independiente con \textsc{Python} y \textsc{Lark}.

\begin{decision}{Alcance}
Construiremos y probaremos únicamente el \textbf{analizador léxico}. La práctica no decidirá todavía si una secuencia de tokens forma una expresión o un comando válido; esa tarea corresponde al análisis sintáctico. Tampoco asignará tipos ni manejará alcances.
\end{decision}

\section*{\cf{verde587246}{Propósito y producto}}

Al terminar, podrás:

\begin{itemize}
  \item explicar por qué el analizador léxico es una fase del compilador y cuál es su contrato;
  \item derivar categorías léxicas a partir de una descripción del lenguaje fuente;
  \item expresar patrones mediante expresiones regulares y relacionarlos con autómatas finitos;
  \item implementar con \textsc{Lark} un lexer independiente para la convención concreta de \textsc{IMP};
  \item conservar atributos y una tabla inicial de localidades;
  \item probar entradas correctas, conflictos entre categorías y caracteres no reconocidos.
\end{itemize}

\begin{evidencia}{Producto esperado}
Una carpeta \texttt{practica1\_analizador\_lexico} que contenga los scripts completados, \texttt{requirements.txt}, ejemplos, pruebas, un \texttt{README.md} y el reporte solicitado para el portafolio. El lexer deberá mostrar por cada token su tipo, lexema, línea y columna.
\end{evidencia}

\section*{\cf{verde587246}{Ruta guiada}}

\begin{center}
\small
\begin{tabular}{|p{0.10\textwidth}|p{0.48\textwidth}|p{0.25\textwidth}|}
\hline
\textbf{Paso} & \textbf{Trabajo principal} & \textbf{Producto parcial} \\
\hline
0 & Preparar el entorno y ejecutar la versión mínima. & Entorno reproducible. \\
\hline
1 & Recuperar el pipeline y el contrato léxico. & Mapa entrada--salida. \\
\hline
2 & Leer la definición de \textsc{IMP} y fijar su escritura concreta. & Inventario de tokens. \\
\hline
3 & Pasar de patrones a expresiones regulares y AFD. & AFD de localidades. \\
\hline
4 & Diseñar reglas, atributos y casos de prueba. & Tabla léxica aprobada. \\
\hline
5 & Extender el lexer mínimo con \textsc{Lark}. & Secuencia de tokens. \\
\hline
6 & Registrar localidades y comprobar posiciones. & Tabla inicial. \\
\hline
7 & Probar errores y preparar la entrega. & Evidencias y reporte. \\
\hline
\end{tabular}
\end{center}

\begin{checkpoint}{Forma de trabajo}
En cada paso: predice el resultado, realiza un cambio pequeño, ejecuta una prueba, compara con otra persona y registra la evidencia. Si el resultado no coincide, explica la diferencia antes de continuar.
\end{checkpoint}

\newpage
\section*{\cf{verde587246}{Paso 0. Preparar el entorno}}

\begin{estacion}{0.1 Copiar el punto de partida}
Copia la carpeta \texttt{scripts\_estudiante} a tu repositorio personal y renómbrala como \texttt{practica1\_analizador\_lexico}. Conserva los archivos y los \texttt{TODO}: indican el orden de trabajo.
\end{estacion}

Estructura inicial:

\begin{lstlisting}[language={}]
practica1_analizador_lexico/
+-- README.md
+-- requirements.txt
+-- afd_localidades.py
+-- lexer_imp.py
+-- probar_lexer.py
+-- pruebas_lexer.py
+-- ejemplos/
    +-- minimo.imp
    +-- flujo.imp
    +-- control.imp
    +-- error.imp
\end{lstlisting}

\begin{estacion}{0.2 Crear un entorno reproducible}
Ejecuta los comandos de \texttt{README.md}. Comprueba que \texttt{requirements.txt} registra una versión concreta de \textsc{Lark}. No instales dependencias adicionales para resolver tareas que pueden realizarse con \textsc{Python}, \textsc{Lark} y \texttt{unittest}.
\end{estacion}

\begin{estacion}{0.3 Ejecutar antes de modificar}
\begin{lstlisting}
python probar_lexer.py ejemplos/minimo.imp
python -m unittest -v pruebas_lexer.py
\end{lstlisting}

Registra:

\begin{itemize}
  \item salida observada: \blank[8cm]
  \item pruebas ejecutadas: \blank[7.6cm]
  \item pruebas omitidas y motivo: \blank[6.1cm]
\end{itemize}
\end{estacion}

\begin{checkpoint}{Checkpoint 0}
\cuadro\ La entrada \texttt{X} produce un token \texttt{LOC}. \quad
\cuadro\ Las pruebas habilitadas terminan correctamente. \quad
\cuadro\ El primer commit conserva el punto de partida.
\end{checkpoint}

\newpage
\section*{\cf{verde587246}{Paso 1. Del software de sistemas al contrato léxico}}

\begin{estacion}{1.1 Recuperar el recorrido}
Completa el esquema con las representaciones involucradas en la transformación.
\end{estacion}

\begin{center}
\begin{tikzpicture}[node distance=1.05cm, every node/.style={font=\small}]
\node[draw,rounded corners,fill=pastelgreen,minimum width=2.4cm,minimum height=.85cm] (f) {programa fuente};
\node[draw,rounded corners,fill=pastelblue,right=of f,minimum width=2.4cm,minimum height=.85cm] (l) {\blank[1.8cm]};
\node[draw,rounded corners,fill=pastelpeach,right=of l,minimum width=2.4cm,minimum height=.85cm] (s) {\blank[1.8cm]};
\node[draw,rounded corners,fill=amber!25,right=of s,minimum width=2.4cm,minimum height=.85cm] (d) {programa destino};
\draw[-{Stealth},thick] (f) -- (l);
\draw[-{Stealth},thick] (l) -- (s);
\draw[-{Stealth},thick] (s) -- (d);
\end{tikzpicture}
\end{center}

\begin{estacion}{1.2 Formular el contrato}
Completa antes de escribir código:

\begin{center}
\begin{tabular}{|p{0.24\textwidth}|p{0.64\textwidth}|}
\hline
\entrada & \blank[9cm] \\[0.45cm]
\hline
\transformacion & Agrupar \blank[3cm], clasificarlos y conservar \blank[3cm]. \\[0.45cm]
\hline
\salida & \blank[9cm] \\[0.45cm]
\hline
Consumidor & \blank[9cm] \\[0.45cm]
\hline
\end{tabular}
\end{center}
\end{estacion}

\begin{modelo}{Ejemplo de una transformación, todavía sin reglas oficiales}
Para el texto \texttt{X := 25}, una salida descriptiva podría separar \texttt{X}, \texttt{:=} y \texttt{25}. La pregunta léxica es qué caracteres forman cada unidad y qué categoría comunica cada una. La pregunta de si la asignación está bien construida se atenderá después.
\end{modelo}

\begin{checkpoint}{Checkpoint 1. Límite de responsabilidad}
Explica por qué encontrar \texttt{X} no permite decidir si fue declarada ni qué valor tiene:

\vspace{1.3cm}
\end{checkpoint}

\newpage
\section*{\cf{verde587246}{Paso 2. Conocer el lenguaje fuente}}

\begin{estacion}{2.1 Sintaxis abstracta de \textsc{IMP}}
Winskel presenta \textsc{IMP} mediante números \(N\), valores de verdad \(T\), localidades \(Loc\), expresiones aritméticas \(Aexp\), expresiones booleanas \(Bexp\) y comandos \(Com\). Sus reglas de formación pueden resumirse así:

\[
\begin{aligned}
a &::= n \mid X \mid a_0+a_1 \mid a_0-a_1 \mid a_0\times a_1,\\
b &::= true \mid false \mid a_0=a_1 \mid a_0\leq a_1
      \mid \neg b \mid b_0\land b_1 \mid b_0\lor b_1,\\
c &::= skip \mid X:=a \mid c_0;c_1
      \mid if\ b\ then\ c_0\ else\ c_1
      \mid while\ b\ do\ c.
\end{aligned}
\]

Estas reglas indican qué construcciones existen, pero no fijan por sí solas todos los detalles de su escritura concreta ni las precedencias.
\end{estacion}

\begin{decision}{Convención concreta del curso}
Para archivos de texto usaremos símbolos ASCII. Esta tabla es una decisión de representación del compilador del curso, no una ampliación de la sintaxis abstracta.
\end{decision}

\begin{center}
\small
\begin{tabular}{|c|c|p{0.48\textwidth}|}
\hline
\textbf{Winskel} & \textbf{Archivo fuente} & \textbf{Decisión} \\
\hline
\(\times\) & \texttt{*} & multiplicación \\
\hline
\(\leq\) & \texttt{<=} & comparación menor o igual \\
\hline
\(\neg\) & \texttt{!} & negación \\
\hline
\(\land\) & \texttt{\&\&} & conjunción \\
\hline
\(\lor\) & \texttt{||} & disyunción \\
\hline
\(n\in N\) & \texttt{NUM} sin signo & el signo menos se emite como token independiente \\
\hline
agrupación concreta & \texttt{(} y \texttt{)} & delimitadores para la gramática posterior \\
\hline
\end{tabular}
\end{center}

\begin{estacion}{2.2 Inventario de tokens}
Completa la columna de atributo. Escribe «no requiere» cuando el nombre del token comunica toda la información necesaria.

\begin{longtable}{|p{0.17\textwidth}|p{0.44\textwidth}|p{0.26\textwidth}|}
\hline
\textbf{Grupo} & \textbf{Tokens y escritura} & \textbf{Atributo} \\
\hline
Datos & \texttt{LOC}: letras seguidas opcionalmente por dígitos & \\
\hline
Datos & \texttt{NUM}: uno o más dígitos & \\
\hline
Valores & \texttt{TRUE}: \texttt{true}; \texttt{FALSE}: \texttt{false} & \\
\hline
Comandos & \texttt{SKIP}, \texttt{IF}, \texttt{THEN}, \texttt{ELSE}, \texttt{WHILE}, \texttt{DO} & \\
\hline
Aritmética & \texttt{PLUS}, \texttt{MINUS}, \texttt{TIMES} & \\
\hline
Booleanos & \texttt{EQ}, \texttt{LE}, \texttt{NOT}, \texttt{AND}, \texttt{OR} & \\
\hline
Control & \texttt{ASSIGN}, \texttt{SEMI}, \texttt{LPAR}, \texttt{RPAR} & \\
\hline
Omitido & espacios, tabuladores y saltos de línea & \\
\hline
\end{longtable}
\end{estacion}

\begin{checkpoint}{Checkpoint 2}
\begin{enumerate}
  \item ¿Por qué \texttt{if} no debe emitirse como \texttt{LOC}? \vspace{0.8cm}
  \item ¿Por qué esta práctica tokeniza \texttt{-} por separado? \vspace{0.8cm}
  \item ¿Qué parte de la tabla deberá respetar la práctica sintáctica? \vspace{0.8cm}
\end{enumerate}
\end{checkpoint}

\newpage
\section*{\cf{verde587246}{Paso 3. De patrones a autómatas}}

\begin{estacion}{3.1 De una descripción a una expresión regular}
Una expresión regular permite especificar con precisión un conjunto de lexemas. Usaremos concatenación, alternativa, clausura \texttt{*}, repetición positiva \texttt{+}, opción \texttt{?} y clases de caracteres.
\end{estacion}

\begin{modelo}{Modelo: numerales}
\[
\text{«uno o más dígitos»}
\quad\longrightarrow\quad
\texttt{[0-9]+}
\]

\begin{center}
\begin{tabular}{|c|c|c|}
\hline
\textbf{Cadena} & \textbf{¿Satisface?} & \textbf{Razón} \\
\hline
\texttt{0} & sí & contiene un dígito \\
\hline
\texttt{25} & sí & contiene dos dígitos \\
\hline
cadena vacía & no & se exige al menos uno \\
\hline
\texttt{-3} & no como un solo \texttt{NUM} & \texttt{-} será otro token \\
\hline
\end{tabular}
\end{center}
\end{modelo}

\begin{estacion}{3.2 Derivar el patrón de localidades}
La descripción de trabajo es: «una o más letras seguidas opcionalmente por cero o más dígitos».

\[
\texttt{LOC} = \blank[7cm]
\]

Prueba tu propuesta con:

\begin{center}
\begin{tabular}{|c|c|p{0.45\textwidth}|}
\hline
\textbf{Cadena} & \textbf{Predicción} & \textbf{Justificación} \\
\hline
\texttt{X} & & \\
\hline
\texttt{contador} & & \\
\hline
\texttt{X2} & & \\
\hline
\texttt{2X} & & \\
\hline
\texttt{X\_2} & & \\
\hline
\texttt{X2Y} & & \\
\hline
\end{tabular}
\end{center}
\end{estacion}

\begin{modelo}{Modelo: AFD de uno o más dígitos}
El estado recuerda si todavía no se ha leído ningún dígito o si ya se leyó al menos uno.

\begin{center}
\begin{tikzpicture}[shorten >=1pt,node distance=3cm,on grid,auto]
\node[state,initial] (q0) {\(q_0\)};
\node[state,accepting,right=of q0] (q1) {\(q_1\)};
\path[->]
  (q0) edge node {\([0-9]\)} (q1)
  (q1) edge[loop above] node {\([0-9]\)} (q1);
\end{tikzpicture}
\end{center}
\end{modelo}

\begin{estacion}{3.3 Convertir y completar el AFD de localidades}
Para \texttt{[A-Za-z]+[0-9]*}, un AFN puede usar una transición \(\varepsilon\) para pasar de la zona de letras a la zona de dígitos. Al construir el AFD, cada estado representa un conjunto de estados alcanzables del AFN. Después se eliminan estados inalcanzables y se comparan estados para decidir si son distinguibles.

Completa la tabla del AFD resultante. Usa \(L\) para letra, \(D\) para dígito y \(O\) para otro carácter.

\begin{center}
\begin{tabular}{|c|c|c|c|c|}
\hline
\textbf{Estado} & \(L\) & \(D\) & \(O\) & \textbf{¿final?} \\
\hline
\(q_0\) & \(q_{\text{letras}}\) & & & no \\
\hline
\(q_{\text{letras}}\) & & & & sí \\
\hline
\(q_{\text{dígitos}}\) & & & & sí \\
\hline
\(q_{\text{error}}\) & & & & no \\
\hline
\end{tabular}
\end{center}

Ahora completa \texttt{TRANSICIONES} en \texttt{afd\_localidades.py} y habilita las pruebas del paso 3.3.
\end{estacion}

\begin{checkpoint}{Checkpoint 3. Conversión y minimización}
\begin{itemize}
  \item \cuadro\ Cada par estado--clase tiene una sola transición.
  \item \cuadro\ \(q_{\text{letras}}\) y \(q_{\text{dígitos}}\) son finales.
  \item \cuadro\ No se fusionan: una letra adicional es válida desde el primero, pero no desde el segundo.
  \item \cuadro\ Las pruebas del AFD producen el resultado predicho.
\end{itemize}
\end{checkpoint}

\newpage
\section*{\cf{verde587246}{Paso 4. Diseñar antes de programar}}

\begin{estacion}{4.1 Tabla de reglas léxicas}
Completa las expresiones o literales de \textsc{Lark}. Usa la tabla como contrato: si cambias un nombre, deberás cambiar las pruebas y documentar la decisión.

\begin{longtable}{|p{0.18\textwidth}|p{0.28\textwidth}|p{0.20\textwidth}|p{0.20\textwidth}|}
\hline
\textbf{Token} & \textbf{Patrón o literal} & \textbf{Acción} & \textbf{Atributo} \\
\hline
\texttt{LOC} & \texttt{/}\blank[2.6cm]\texttt{/} & emitir y registrar & \\
\hline
\texttt{NUM} & \texttt{/}\blank[2.6cm]\texttt{/} & emitir & \\
\hline
\texttt{IF} & \texttt{"if"} & emitir & \\
\hline
\texttt{THEN} & & emitir & \\
\hline
\texttt{ELSE} & & emitir & \\
\hline
\texttt{WHILE} & & emitir & \\
\hline
\texttt{DO} & & emitir & \\
\hline
\texttt{SKIP} & & emitir & \\
\hline
\texttt{TRUE}, \texttt{FALSE} & & emitir & \\
\hline
\texttt{ASSIGN} & & emitir & \\
\hline
\texttt{LE} & & emitir & \\
\hline
\texttt{PLUS}, \texttt{MINUS}, \texttt{TIMES} & & emitir & \\
\hline
\texttt{EQ}, \texttt{NOT}, \texttt{AND}, \texttt{OR} & & emitir & \\
\hline
\texttt{SEMI}, \texttt{LPAR}, \texttt{RPAR} & & emitir & \\
\hline
\texttt{WS} & clase importada de \textsc{Lark} & omitir & no aplica \\
\hline
\end{longtable}
\end{estacion}

\begin{estacion}{4.2 Resolver una colisión}
El lexema \texttt{if} satisface también el patrón de \texttt{LOC}. El lexer básico no consulta al parser para decidir. En \textsc{Lark} combinaremos una prioridad explícita con un límite léxico:

\begin{lstlisting}[language={}]
IF.2: /if(?![A-Za-z0-9])/
LOC: /[A-Za-z]+[0-9]*/
\end{lstlisting}

La prioridad resuelve el empate cuando la entrada es exactamente \texttt{if}; la anticipación negativa \texttt{(?![A-Za-z0-9])} impide que la palabra reservada corte una localidad más larga. Así, \texttt{if2} e \texttt{ifX} pueden reconocerse completas como \texttt{LOC}.

Predice:

\begin{center}
\begin{tabular}{|c|c|c|}
\hline
\textbf{Entrada} & \textbf{Token esperado} & \textbf{Razón} \\
\hline
\texttt{if} & & \\
\hline
\texttt{if2} & & \\
\hline
\texttt{ifX} & & \\
\hline
\end{tabular}
\end{center}
\end{estacion}

\begin{estacion}{4.3 Especificar pruebas antes del cambio}
Escribe el resultado esperado para:

\begin{lstlisting}[language={}]
X := (Y + 2) <= 10 && true;
\end{lstlisting}

\vspace{1.5cm}

Compara tu lista con la prueba \texttt{test\_operadores\_y\_delimitadores}, pero no retires todavía su decorador \texttt{@unittest.skip}.
\end{estacion}

\begin{checkpoint}{Checkpoint 4}
La tabla debe permitir responder, para cada categoría: qué reconoce, qué emite, qué conserva y con qué prueba se verificará.
\end{checkpoint}

\newpage
\section*{\cf{verde587246}{Paso 5. Implementar con \textsc{Lark}}}

\begin{modelo}{5.1 Versión mínima}
El archivo inicial contiene:

\begin{lstlisting}
from lark import Lark

GRAMATICA_LEXICA = r"""
LOC: /[A-Za-z]+[0-9]*/
%import common.WS
%ignore WS
"""

lexer = Lark(
    GRAMATICA_LEXICA,
    parser=None,
    lexer="basic",
)

tokens = list(lexer.lex("X"))
\end{lstlisting}

\texttt{parser=None} mantiene el trabajo en la fase léxica; \texttt{lexer="basic"} selecciona un lexer independiente; \texttt{lex()} recorre la entrada y produce tokens.
\end{modelo}

\begin{estacion}{5.2 Incorporar numerales}
\begin{enumerate}
  \item Agrega \texttt{NUM} a \texttt{GRAMATICA\_LEXICA}.
  \item Ejecuta \texttt{python probar\_lexer.py ejemplos/minimo.imp}.
  \item Habilita \texttt{test\_localidades\_y\_numerales}.
  \item Ejecuta las pruebas y registra el resultado.
\end{enumerate}

Commit o registro del cambio: \blank[8cm]
\end{estacion}

\begin{estacion}{5.3 Incorporar palabras reservadas}
\begin{enumerate}
  \item Agrega las ocho palabras reservadas con prioridad superior a \texttt{LOC} y con el límite léxico del modelo de \texttt{IF}.
  \item Habilita \texttt{test\_palabras\_reservadas\_no\_son\_localidades}.
  \item Añade un caso propio con una localidad que empiece igual que una palabra reservada.
\end{enumerate}

Caso propio y resultado: \blank[8cm]
\end{estacion}

\begin{estacion}{5.4 Incorporar símbolos}
Agrega operadores, asignación, secuenciación y paréntesis. Después:

\begin{lstlisting}
python -m unittest -v pruebas_lexer.py
python probar_lexer.py ejemplos/flujo.imp
python probar_lexer.py ejemplos/control.imp
\end{lstlisting}

Habilita \texttt{test\_operadores\_y\_delimitadores}. No modifiques la salida esperada para hacer pasar una implementación incorrecta; compara primero con la tabla aprobada.
\end{estacion}

\begin{evidencia}{Evidencia 5}
Conserva una captura o transcripción legible de:
\begin{itemize}
  \item la ejecución de las pruebas habilitadas;
  \item la secuencia producida para \texttt{flujo.imp};
  \item la secuencia producida para \texttt{control.imp};
  \item el commit o registro que corresponde a cada incremento.
\end{itemize}
\end{evidencia}

\newpage
\section*{\cf{verde587246}{Paso 6. Atributos y tabla inicial de localidades}}

\begin{estacion}{6.1 Observar atributos}
Cada token de \textsc{Lark} conserva su tipo, lexema y posición. Ejecuta:

\begin{lstlisting}
python probar_lexer.py ejemplos/flujo.imp
\end{lstlisting}

Selecciona dos tokens del mismo tipo con lexemas distintos:

\begin{center}
\begin{tabular}{|c|c|c|c|}
\hline
\textbf{Tipo} & \textbf{Lexema} & \textbf{Línea} & \textbf{Columna} \\
\hline
 & & & \\
\hline
 & & & \\
\hline
\end{tabular}
\end{center}

¿Qué información permanece igual y cuál cambia? \vspace{1cm}
\end{estacion}

\begin{estacion}{6.2 Construir la tabla}
Completa \texttt{construir\_tabla\_localidades} en \texttt{lexer\_imp.py}. La tabla de esta práctica sólo registra:

\begin{itemize}
  \item lexema de la localidad;
  \item línea y columna de la primera aparición;
  \item número total de apariciones.
\end{itemize}

No agregues tipos, valores ni alcances: aún no contamos con información semántica para hacerlo.

Habilita \texttt{test\_tabla\_de\_localidades} y ejecuta:

\begin{lstlisting}
python probar_lexer.py ejemplos/flujo.imp --tabla
\end{lstlisting}
\end{estacion}

\begin{checkpoint}{Checkpoint 6. Separar responsabilidades}
Completa:

\begin{itemize}
  \item El lexer puede registrar que apareció \texttt{X} en \blank[3cm].
  \item Todavía no puede afirmar que \texttt{X} fue \blank[4cm].
  \item El tipo de \texttt{X} se comprobará en la fase de \blank[4cm].
\end{itemize}
\end{checkpoint}

\newpage
\section*{\cf{verde587246}{Paso 7. Errores, pruebas y entrega}}

\begin{estacion}{7.1 Diagnóstico léxico}
Ejecuta:

\begin{lstlisting}
python probar_lexer.py ejemplos/error.imp
\end{lstlisting}

El carácter \texttt{@} no inicia ningún patrón. El controlador captura \texttt{UnexpectedCharacters} e informa línea y columna. Esta práctica se detiene en el primer error: no inventaremos una corrección ni una estrategia de recuperación.
\end{estacion}

\begin{estacion}{7.2 Activar y ampliar pruebas}
Habilita \texttt{test\_caracter\_desconocido}. Después agrega, como mínimo, un caso propio de cada clase:

\begin{center}
\begin{tabular}{|p{0.30\textwidth}|p{0.28\textwidth}|p{0.28\textwidth}|}
\hline
\textbf{Clase} & \textbf{Entrada} & \textbf{Resultado esperado} \\
\hline
Mínimo correcto & & \\
\hline
Varias categorías & & \\
\hline
Palabra reservada frente a \texttt{LOC} & & \\
\hline
Operador de dos caracteres & & \\
\hline
Espacios y saltos de línea & & \\
\hline
Carácter desconocido & & \\
\hline
Caso frontera propio & & \\
\hline
\end{tabular}
\end{center}
\end{estacion}

\begin{checkpoint}{Checkpoint 7. Prueba completa}
\begin{itemize}
  \item \cuadro\ Todas las pruebas habilitadas terminan correctamente.
  \item \cuadro\ Cada prueba tiene un propósito identificable.
  \item \cuadro\ Los errores informan una posición relacionada con la entrada.
  \item \cuadro\ No se usa el parser para resolver decisiones léxicas.
  \item \cuadro\ No se añadieron comentarios u otras construcciones no aprobadas para \textsc{IMP}.
\end{itemize}
\end{checkpoint}

\section*{\cf{verde587246}{Criterios de entrega}}

\begin{longtable}{|p{0.32\textwidth}|p{0.55\textwidth}|}
\hline
\textbf{Criterio} & \textbf{Evidencia observable} \\
\hline
Especificación & Tabla de tokens consistente con \textsc{IMP} y con la convención concreta. \\
\hline
Fundamento formal & Expresiones regulares justificadas y AFD de localidades funcional. \\
\hline
Implementación & Lexer independiente con \textsc{Lark}; reglas legibles y nombres consistentes. \\
\hline
Atributos & Tipo, lexema y posición visibles; tabla inicial de localidades correcta. \\
\hline
Pruebas & Casos correctos, colisiones, espacios, operadores compuestos y error léxico. \\
\hline
Documentación & \texttt{README.md}, instalación, uso, decisiones y limitaciones. \\
\hline
Evolución & Historial de \textsc{Git} que permita reconocer incrementos y verificaciones. \\
\hline
Reporte & Introducción, desarrollo, evidencias, resultados, errores, conclusiones, tiempos y vínculo al repositorio. \\
\hline
\end{longtable}

\begin{evidencia}{Lista final}
\begin{itemize}
  \item \cuadro\ \texttt{requirements.txt} fija la versión de \textsc{Lark}.
  \item \cuadro\ Los scripts no conservan \texttt{TODO} pendientes de la práctica.
  \item \cuadro\ Las pruebas ya no están omitidas.
  \item \cuadro\ Los ejemplos pueden ejecutarse desde las instrucciones del \texttt{README.md}.
  \item \cuadro\ El reporte distingue lo construido por el equipo de lo automatizado por \textsc{Lark}.
  \item \cuadro\ La salida léxica queda preparada para ser consumida por la práctica sintáctica.
\end{itemize}
\end{evidencia}

\section*{\cf{verde587246}{Bitácora y cierre}}

\begin{center}
\begin{tabular}{|p{0.32\textwidth}|p{0.20\textwidth}|p{0.20\textwidth}|}
\hline
\textbf{Etapa} & \textbf{Estimación} & \textbf{Tiempo real} \\
\hline
Preparación y recuperación & & \\
\hline
Expresiones regulares y AFD & & \\
\hline
Diseño de reglas & & \\
\hline
Implementación incremental & & \\
\hline
Tabla, pruebas y documentación & & \\
\hline
\textbf{Total} & & \\
\hline
\end{tabular}
\end{center}

\begin{checkpoint}{Reflexión final para el reporte}
\begin{enumerate}
  \item ¿Qué decisión léxica resultó más difícil y con qué evidencia la resolviste?
  \item ¿Qué diferencia observaste entre describir un patrón, reconocerlo con un AFD e implementarlo con \textsc{Lark}?
  \item ¿Qué información deberá conservarse sin cambios cuando se integre el analizador sintáctico?
\end{enumerate}
\end{checkpoint}

\section*{\cf{verde587246}{Conclusión}}

El analizador construido materializa la primera transformación interna del compilador: recibe caracteres, reconoce lexemas, emite tokens y conserva atributos verificables. Las expresiones regulares especifican las categorías; los autómatas explican su reconocimiento; y \textsc{Lark} automatiza parte de la ejecución, pero no decide el vocabulario, los nombres, las prioridades ni las pruebas. El contrato documentado permitirá integrar esta salida con el analizador sintáctico en una práctica posterior.

\section*{\cf{verde587246}{Referencias}}

\begin{enumerate}[label=\textbf{[\arabic*]}]
  \item Winskel, G. (1993). \textit{The Formal Semantics of Programming Languages: An Introduction}. The MIT Press.
  \item Aho, A. V., Lam, M. S., Sethi, R., y Ullman, J. D. (2008). \textit{Compiladores: principios, técnicas y herramientas} (2.\textsuperscript{a} ed.). Pearson Educación. ISBN 978-970-26-1133-2.
  \item Thain, D. (2026). \textit{Introduction to Compilers and Language Design} (2.\textsuperscript{a} ed.). University of Notre Dame.
  \item Vilar Torres, J. M. (2010). \textit{Estructura de los compiladores e intérpretes}. Universitat Jaume I.
  \item Vilar Torres, J. M. (2010). \textit{Analizador léxico}. Universitat Jaume I.
  \item Lark. (s.\,f.). \textit{API Reference y Grammar Reference}. \url{https://lark-parser.readthedocs.io/}
\end{enumerate}

\end{document}
